\chapter{Introduction}
\label{ch:intro}

% TODO: Replace placeholders with actual content. Invoke /draft-chapter to generate a first draft.

% Outline approved by outline-curator agent. Sections mirror course/lectures/01-intro/notes.md.

% ---------------------------------------------------------------
\section{A Brief History of Textual Analysis in Finance}
\label{sec:history}

\subsection{Early Approaches: Keyword Counting and Dictionary Methods}
\label{sec:history-early}

[Placeholder]

\subsection{The Shift to Statistical and Neural Methods}
\label{sec:history-shift}

[Placeholder]

% ---------------------------------------------------------------
\section{Why Text Matters in Finance}
\label{sec:motivation}

\subsection{The Information Content of Earnings Calls, News, and Filings}
\label{sec:motivation-info}

[Placeholder]

\subsection{Text as a Signal: Evidence from the Literature}
\label{sec:motivation-evidence}

[Placeholder]

% ---------------------------------------------------------------
\section{Classical Text Representations}
\label{sec:classical}

\subsection{Vocabulary, Tokens, and the Bag-of-Words Model}
\label{sec:classical-bow}

[Placeholder]

\subsection{One-Hot Encoding}
\label{sec:classical-onehot}

[Placeholder]

\subsection{TF-IDF Weighting}
\label{sec:classical-tfidf}

[Placeholder]

% ---------------------------------------------------------------
\section{Word Embeddings}
\label{sec:embeddings}

The Bag-of-Words and TF-IDF representations introduced in
Section~\ref{sec:classical} treat every word as an isolated symbol.
Two terms that never co-occur with each other but share essentially
the same contextual neighbours — \emph{equity} and \emph{stock}, for
instance — receive orthogonal vectors with cosine similarity exactly
zero.  This is not merely aesthetically unsatisfying: it means that a
model trained on documents containing the word \emph{equity} will
generalise poorly to documents using the synonym \emph{stock}, even
though the two words carry the same information.  Word embeddings solve
this problem by mapping every vocabulary item to a dense, low-dimensional
real-valued vector so that semantically related words land near one
another in the embedding space.

% ----------------------------------------------------------------
\subsection{Distributional Semantics and the Embedding Idea}
\label{sec:embeddings-idea}

\paragraph{The distributional hypothesis.}
The theoretical foundation of word embeddings is the
\emph{distributional hypothesis}, articulated most famously by the
British linguist J.R.\ Firth:
\begin{quote}
  ``You shall know a word by the company it keeps.''
  \hfill\cite{firth1957} % [CITE: Firth, J.R. (1957). Papers in Linguistics 1934–1951]
\end{quote}
The idea is that the meaning of a word can be inferred from the
distribution of other words that appear in its vicinity across a large
corpus.  Words with similar meanings tend to appear in similar
contexts; therefore a representation that captures contextual
co-occurrence patterns will capture semantic similarity.

\paragraph{One-hot encoding and its sparsity problem.}
Recall that the classical one-hot encoding assigns to each word
$w \in \mathcal{V}$ the indicator vector $\mathbf{e}_w \in \{0,1\}^{|\mathcal{V}|}$
with a single non-zero entry at the position indexed by $w$.  This
representation has two critical deficiencies:
\begin{enumerate}
  \item \textbf{Sparsity.} For a vocabulary of $|\mathcal{V}| = 50{,}000$ tokens,
        each word is represented by a vector that is $50{,}000$-dimensional
        but contains exactly one non-zero entry.  Storage, computation, and
        learning all scale poorly with this dimensionality.
  \item \textbf{No semantic structure.} Every pair of one-hot vectors is
        orthogonal: $\mathbf{e}_i^\top \mathbf{e}_j = \mathbf{1}[i = j]$.
        The representation encodes no information about semantic relatedness.
\end{enumerate}

\paragraph{Word embeddings: formal definition.}
A word embedding is a function that maps discrete vocabulary items to
dense, low-dimensional vectors.

\begin{definition}[Word Embedding]
\label{def:word-embedding}
Let $\mathcal{V}$ be a finite vocabulary of size $V = |\mathcal{V}|$.
A \emph{word embedding} of dimension $d$ is an injective map
\begin{equation}
  \varphi : \mathcal{V} \;\longrightarrow\; \mathbb{R}^{d},
  \qquad d \ll V.
  \label{eq:embedding-map}
\end{equation}
The matrix $\mathbf{W} \in \mathbb{R}^{V \times d}$ whose $i$-th row is
$\varphi(w_i)$ is called the \emph{embedding matrix}.  Retrieving the
embedding of word $w_i$ is equivalent to the matrix--vector product
$\mathbf{W}^\top \mathbf{e}_{w_i}$, i.e.\ selecting the $i$-th row of
$\mathbf{W}$.
\end{definition}

The embedding dimension $d$ is a hyperparameter.  Typical values in
practice range from $d = 50$ (GloVe-50) to $d = 300$ (GloVe-300,
Word2Vec) and up to $d = 768$ or $d = 1024$ for transformer-based
contextual embeddings (see Section~\ref{sec:doc-repr-sbert}).  The key
constraint is $d \ll V$: whereas one-hot encoding requires $V$ numbers
per word, an embedding requires only $d$ numbers, a compression ratio
of $V/d \approx 150$--$1000$ for typical values.

\begin{remark}
\label{rem:embedding-geometry}
The power of word embeddings lies not only in their compactness but in
their geometry.  Under a well-trained embedding, the Euclidean distance
or cosine similarity between two embedding vectors reflects the semantic
distance between the corresponding words.  Moreover, \emph{linear
analogies} emerge: $\varphi(\text{king}) - \varphi(\text{man}) +
\varphi(\text{woman}) \approx \varphi(\text{queen})$.  In a financial
corpus one observes analogies such as $\varphi(\text{call option}) -
\varphi(\text{upside}) + \varphi(\text{downside}) \approx
\varphi(\text{put option})$.
\end{remark}

% ----------------------------------------------------------------
\subsection{Word2Vec and GloVe}
\label{sec:embeddings-word2vec}

Two landmark algorithms — Word2Vec \cite{mikolov2013} % [CITE: Mikolov et al. (2013). Distributed Representations of Words and Phrases and their Compositionality. NeurIPS]
and GloVe \cite{pennington2014} % [CITE: Pennington, Socher, Manning (2014). GloVe: Global Vectors for Word Representation. EMNLP]
— remain the canonical dense embedding methods.  We treat each in turn,
then discuss finance-specific adaptations.

% ............................................................
\subsubsection{Word2Vec: Skip-Gram with Negative Sampling}

Word2Vec frames the embedding problem as a supervised prediction task.
The \emph{Skip-Gram} variant trains a word to predict the words that
appear in a sliding window of $\pm c$ positions around it.

\paragraph{Notation.}
Let the corpus consist of a sequence of $T$ tokens
$w_1, w_2, \dots, w_T$ drawn from vocabulary $\mathcal{V}$.  Each
word $w \in \mathcal{V}$ is associated with two embedding vectors:
a \emph{centre-word} vector $\hat{\mathbf{v}}_w \in \mathbb{R}^d$ (used when
$w$ is the target word) and a \emph{context-word} vector
$\mathbf{v}_w \in \mathbb{R}^d$ (used when $w$ is a context word).
This two-vector convention simplifies the mathematics and is discarded
after training (often the two sets are averaged or the context vectors
are thrown away).

\paragraph{The Skip-Gram objective.}
For each centre word $w_t$ and each context position $k \in
\{-c, \ldots, -1, 1, \ldots, c\}$, the model defines the conditional
probability of observing context word $w_{t+k}$ given centre word
$w_t$ via a soft-max over the full vocabulary:
\begin{equation}
  P(w_{t+k} \mid w_t)
  \;=\;
  \frac{\exp\!\bigl(\mathbf{v}_{w_{t+k}}^{\top} \hat{\mathbf{v}}_{w_t}\bigr)}
       {\displaystyle\sum_{w=1}^{V} \exp\!\bigl(\mathbf{v}_{w}^{\top} \hat{\mathbf{v}}_{w_t}\bigr)}.
  \label{eq:skipgram-softmax}
\end{equation}
The Skip-Gram objective maximises the average log-likelihood of all
centre--context pairs:
\begin{equation}
  J_{\mathrm{SG}}
  \;=\;
  \frac{1}{T}
  \sum_{t=1}^{T}
  \sum_{\substack{k=-c \\ k \neq 0}}^{c}
  \log P(w_{t+k} \mid w_t).
  \label{eq:skipgram-obj}
\end{equation}

\paragraph{Computational intractability and negative sampling.}
The normalising constant in \eqref{eq:skipgram-softmax} sums over
all $V \approx 10^4$--$10^5$ words at every gradient step, making
exact optimisation prohibitively expensive.  \emph{Negative Sampling}
(NEG) replaces the full soft-max with a binary classification task:
distinguish the observed context word (positive sample) from $K$
words drawn from a noise distribution $P_n$ (negative samples).  The
NEG objective for a single centre--context pair $(w_t, w_{t+k})$ is
\begin{equation}
  \mathcal{L}_{\mathrm{NEG}}(w_t, w_{t+k})
  \;=\;
  \log \sigma\!\bigl(\mathbf{v}_{w_{t+k}}^{\top} \hat{\mathbf{v}}_{w_t}\bigr)
  \;+\;
  \sum_{j=1}^{K}
  \mathbb{E}_{w_j \sim P_n}\!\left[
    \log \sigma\!\bigl(-\mathbf{v}_{w_j}^{\top} \hat{\mathbf{v}}_{w_t}\bigr)
  \right],
  \label{eq:neg-sampling-obj}
\end{equation}
where $\sigma(x) = (1 + e^{-x})^{-1}$ is the logistic sigmoid.
Mikolov et al.\ recommend $K = 5$--$20$ for small corpora and $K = 2$--$5$
for large corpora, with $P_n(w) \propto \text{freq}(w)^{3/4}$ (the
unigram distribution raised to the $3/4$ power), which down-weights
the most frequent stop-words.

\paragraph{Gradient update.}
Taking the gradient of \eqref{eq:neg-sampling-obj} with respect to the
context vector $\mathbf{v}_{w_{t+k}}$ (positive sample) gives
\begin{equation}
  \frac{\partial \mathcal{L}_{\mathrm{NEG}}}{\partial \mathbf{v}_{w_{t+k}}}
  \;=\;
  \bigl[1 - \sigma(\mathbf{v}_{w_{t+k}}^{\top} \hat{\mathbf{v}}_{w_t})\bigr]
  \hat{\mathbf{v}}_{w_t}
  \;=\;
  \bigl[1 - P_{\mathrm{NEG}}(D=1 \mid w_{t+k}, w_t)\bigr]\hat{\mathbf{v}}_{w_t},
  \label{eq:neg-grad-pos}
\end{equation}
and for each negative sample $w_j$:
\begin{equation}
  \frac{\partial \mathcal{L}_{\mathrm{NEG}}}{\partial \mathbf{v}_{w_j}}
  \;=\;
  -\sigma(\mathbf{v}_{w_j}^{\top} \hat{\mathbf{v}}_{w_t})\,\hat{\mathbf{v}}_{w_t}.
  \label{eq:neg-grad-neg}
\end{equation}
The corresponding update for the centre vector $\hat{\mathbf{v}}_{w_t}$ is
obtained by symmetry.  These updates have the appealing interpretation
that the positive context vector is pushed closer to the centre vector
(its coefficient $1 - \sigma(\cdot)$ is large when the two are
dissimilar), while each negative sample's vector is pushed away.
Stochastic gradient descent with these updates is extremely efficient:
each step touches only $K + 1$ rows of the embedding matrix rather
than all $V$ rows.

\paragraph{Analogy structure.}
A celebrated property of Word2Vec embeddings is the emergence of
\emph{linear analogical relationships}.  The classic example is
\begin{equation}
  \hat{\mathbf{v}}_{\text{king}}
  - \hat{\mathbf{v}}_{\text{man}}
  + \hat{\mathbf{v}}_{\text{woman}}
  \;\approx\;
  \hat{\mathbf{v}}_{\text{queen}}.
  \label{eq:analogy-royalty}
\end{equation}
In a financial corpus one encounters analogies such as
$\hat{\mathbf{v}}_{\text{equity}} - \hat{\mathbf{v}}_{\text{dividend}}
+ \hat{\mathbf{v}}_{\text{coupon}} \approx \hat{\mathbf{v}}_{\text{bond}}$.
These regularities arise because the embedding must be consistent with
patterns of co-occurrence: if ``king'' and ``queen'' share the same
contextual environments except for gender-related words, their
difference captures a gender direction.

% ............................................................
\subsubsection{GloVe: Global Vectors for Word Representation}

Whereas Word2Vec processes one context window at a time,
\emph{Global Vectors} (GloVe) \cite{pennington2014} % [CITE: Pennington et al. (2014) EMNLP]
directly factorises the global word--word co-occurrence matrix,
explicitly encoding corpus-level statistics.

\paragraph{Co-occurrence matrix.}
Let $X_{ij}$ denote the number of times word $j$ appears in the
context window of word $i$ across the entire corpus, and let
$X_i = \sum_j X_{ij}$ be the total context count for word $i$.

\begin{definition}[Co-occurrence Matrix]
\label{def:cooccurrence}
Given a corpus of $T$ tokens, a vocabulary $\mathcal{V}$ of size $V$,
and a context window of half-width $c$, the \emph{co-occurrence matrix}
$\mathbf{X} \in \mathbb{N}^{V \times V}$ has entries
\begin{equation}
  X_{ij}
  \;=\;
  \sum_{t=1}^{T}
  \sum_{\substack{k=-c \\ k \neq 0}}^{c}
  \mathbf{1}[w_t = w_i]\,\mathbf{1}[w_{t+k} = w_j].
  \label{eq:cooccurrence}
\end{equation}
\end{definition}

\paragraph{The GloVe objective.}
The key theoretical insight of GloVe is that the \emph{ratio} of
co-occurrence probabilities $P_{ik}/P_{jk}$ encodes relational
meaning more faithfully than the raw probabilities.  This motivates
learning vectors $\mathbf{w}_i, \tilde{\mathbf{w}}_j \in \mathbb{R}^d$
and scalar biases $b_i, \tilde{b}_j \in \mathbb{R}$ by minimising the
weighted squared loss:
\begin{equation}
  J_{\mathrm{GloVe}}
  \;=\;
  \sum_{i,j=1}^{V}
  f(X_{ij})
  \Bigl(
    \mathbf{w}_i^{\top} \tilde{\mathbf{w}}_j
    + b_i + \tilde{b}_j
    - \log X_{ij}
  \Bigr)^{2},
  \label{eq:glove-obj}
\end{equation}
where $f : \mathbb{R}_{\geq 0} \to \mathbb{R}_{\geq 0}$ is a weighting
function that down-weights both very rare and very frequent co-occurrences:
\begin{equation}
  f(x)
  \;=\;
  \begin{cases}
    \bigl(x / x_{\max}\bigr)^{\alpha} & \text{if } x < x_{\max}, \\
    1 & \text{otherwise.}
  \end{cases}
  \label{eq:glove-weight}
\end{equation}
Pennington et al.\ recommend $\alpha = 0.75$ and $x_{\max} = 100$.
Pairs with $X_{ij} = 0$ contribute nothing to the sum (since
$f(0) = 0$); the logarithm is therefore only evaluated at positive
counts.

\begin{remark}
\label{rem:glove-vs-word2vec}
GloVe and Word2Vec learn similar representations in practice, but their
optimisation philosophies differ.  GloVe is a single-pass weighted
least-squares problem over the co-occurrence matrix, which is assembled
once and stored; Word2Vec performs online stochastic gradient descent
over individual windows without ever materialising the full matrix.  For
large vocabularies, storing $\mathbf{X}$ requires $O(V^2)$ memory in
the dense case, though in practice $\mathbf{X}$ is extremely sparse.
GloVe tends to be faster to train when co-occurrence counts can be
computed in a preprocessing step, while Word2Vec scales more gracefully
to streaming corpora.
\end{remark}

% ............................................................
\subsubsection{Finance-Specific Embeddings and Domain Adaptation}

General-purpose embeddings trained on Wikipedia or Common Crawl encode
the general-English meaning of financial terms, which can differ
substantially from their technical meanings.  The word \emph{risk}, for
example, has nearest neighbours \{danger, threat, hazard, peril\} in a
general embedding but nearest neighbours \{volatility, variance,
exposure, drawdown, VaR\} in a finance-domain embedding.  Similarly,
\emph{short} means ``brief'' in general English but ``short-sell a
security'' in a trading context.

Domain adaptation strategies include:
\begin{enumerate}
  \item \textbf{Domain-corpus pre-training.}  Train Word2Vec or GloVe
        from scratch on a large financial corpus (SEC filings, Bloomberg
        articles, earnings call transcripts).  The Bloomberg Financial
        Word2Vec model \cite{liu2018} % [CITE: Liu et al. (2018). Efficient low-rank multimodal fusion with modality-specific factors. ACL — or appropriate FinNLP reference]
        and the FinancialPhraseBank embeddings are examples.
  \item \textbf{Fine-tuned transformer embeddings.}  FinBERT
        \cite{yang2020} % [CITE: Yang et al. (2020). FinBERT: A Pretrained Language Model for Financial Communications. arXiv:2006.08097]
        is a BERT model \cite{devlin2019} % [CITE: Devlin et al. (2019). BERT: Pre-training of Deep Bidirectional Transformers for Language Understanding. NAACL]
        further pre-trained on Reuters news, 10-K filings, and earnings
        calls.  Its token embeddings capture the financial sense of
        polysemous terms such as \emph{margin}, \emph{return}, and
        \emph{position}.
  \item \textbf{Retrofitting.}  Post-hoc adjustment of a general
        embedding using a domain-specific lexicon or knowledge graph
        (e.g.\ a financial ontology) to pull synonymous terms together
        and push antonyms apart.
\end{enumerate}

\begin{example}
\label{ex:risk-neighbours}
Table~\ref{tab:risk-neighbours} illustrates the five nearest neighbours
of the word \emph{risk} under cosine similarity in (a) a GloVe model
trained on Wikipedia and (b) a GloVe model trained on a corpus of 1M
SEC annual reports (10-K filings).

\begin{table}[H]
  \centering
  \caption{Nearest neighbours of \emph{risk} in two embedding spaces.
           Finance-domain training shifts the semantic neighbourhood
           from colloquial to technical usage.}
  \label{tab:risk-neighbours}
  \begin{tabular}{llll}
    \toprule
    Rank & General (Wikipedia) & Finance (10-K) \\
    \midrule
    1 & danger      & uncertainty   \\
    2 & threat      & exposure      \\
    3 & hazard      & volatility    \\
    4 & possibility & credit risk   \\
    5 & factor      & downside      \\
    \bottomrule
  \end{tabular}
\end{table}
\end{example}

% ----------------------------------------------------------------
\subsection{Case Study: PCA on Financial Vocabulary}
\label{sec:embeddings-pca}

This case study provides a concrete, visual demonstration that word
embeddings encode domain-specific semantic structure.  The analysis is
implemented in full in
\texttt{code/notebooks/01-intro/demo.ipynb}.

\paragraph{Setup.}
We load a pre-trained GloVe model (either the 300-dimensional
\texttt{glove.6B.300d} release or a finance-tuned variant) and extract
the embedding vectors for the following 30 financial terms:
\begin{center}
  \textit{equity, bond, yield, spread, volatility, risk, return,
  portfolio, hedge, derivative, option, futures, dividend, earnings,
  revenue, profit, loss, merger, acquisition, default, credit,
  inflation, central bank, interest rate, liquidity, leverage,
  sentiment, alpha, beta, factor.}
\end{center}
These terms span four semantic groups:
\begin{enumerate}
  \item \textbf{Rates \& macro:} yield, spread, inflation, central bank,
        interest rate.
  \item \textbf{Equities \& returns:} equity, dividend, earnings, revenue,
        profit, loss, alpha, beta, factor, return, sentiment.
  \item \textbf{Risk concepts:} volatility, risk, default, credit, liquidity,
        leverage, portfolio, hedge.
  \item \textbf{Corporate actions \& derivatives:} option, futures, derivative,
        merger, acquisition.
\end{enumerate}

\paragraph{Dimensionality reduction.}
Let $\mathbf{E} \in \mathbb{R}^{30 \times d}$ be the matrix whose
$i$-th row is the (mean-centred) embedding vector of the $i$-th term.
We apply principal component analysis (PCA) to obtain a
$30 \times 2$ matrix of scores:
\begin{equation}
  \mathbf{Z}
  \;=\;
  \mathbf{E}\,\mathbf{P}_{:,1:2},
  \label{eq:pca-projection}
\end{equation}
where $\mathbf{P} \in \mathbb{R}^{d \times d}$ contains the principal
components (eigenvectors of $\mathbf{E}^{\top}\mathbf{E}$) sorted by
descending eigenvalue, and $\mathbf{P}_{:,1:2}$ retains the first two.

\paragraph{Results and interpretation.}
Figure~\ref{fig:embeddings-pca} shows the scatter plot of the 30 terms
in the two-dimensional PCA space, coloured by semantic group.

\begin{figure}[H]
  \centering
  \includegraphics[width=0.82\textwidth]{figures/embeddings-pca}
  \caption{PCA projection of embedding vectors for 30 financial terms
           into two dimensions.  Points are coloured by semantic group:
           \textcolor{blue}{rates \& macro} (blue),
           \textcolor{green!60!black}{equities \& returns} (green),
           \textcolor{red}{risk concepts} (red), and
           \textcolor{orange}{corporate actions \& derivatives} (orange).
           Clusters reveal that the embedding space encodes domain-specific
           semantic proximity: \emph{option} and \emph{futures} lie close
           to \emph{derivative}; \emph{default} and \emph{credit} co-cluster;
           and the rates cluster (\emph{yield}, \emph{spread},
           \emph{interest rate}) is well-separated from the equities cluster.}
  \label{fig:embeddings-pca}
\end{figure}

Several patterns are consistently observable across different
pre-trained models:
\begin{itemize}
  \item \emph{option}, \emph{futures}, and \emph{derivative} form a tight
        cluster, reflecting their shared contextual vocabulary (premium,
        contract, expiry, strike).
  \item \emph{default} and \emph{credit} lie close together and at a
        distance from \emph{equity}-group terms, consistent with the
        credit-versus-equity divide in fixed income.
  \item \emph{merger} and \emph{acquisition} are nearly synonymous in
        this corpus, mapping to almost identical embedding vectors.
  \item \emph{alpha}, \emph{beta}, and \emph{factor} cluster together,
        reflecting their shared origin in the asset-pricing literature
        and their frequent co-occurrence in portfolio management texts.
  \item The rates cluster (\emph{yield}, \emph{spread}, \emph{inflation},
        \emph{interest rate}, \emph{central bank}) is the most cohesive,
        owing to the highly formulaic language of fixed-income analysis.
\end{itemize}

\begin{remark}
\label{rem:pca-directions}
The principal component directions $\mathbf{p}_1, \mathbf{p}_2$
are not directly interpretable as named financial concepts.  They are
the linear combinations of the original $d = 300$ embedding dimensions
that explain the most variance \emph{among the 30 selected terms}.  A
different selection of terms would produce different axes.  The visual
clustering is informative, but one should resist assigning
domain-specific names to the $x$- and $y$-axes without rigorous
validation.  Rotation methods (e.g.\ Varimax) and supervised probing
classifiers offer more principled approaches to identifying
interpretable directions in embedding space; see \cite{tenney2019} % [CITE: Tenney et al. (2019). BERT Rediscovers the Classical NLP Pipeline. ACL]
for an overview in the transformer setting.
\end{remark}

% ---------------------------------------------------------------
\section{Document-Level Representations}
\label{sec:doc-repr}

Word embeddings assign vectors to individual tokens.  In most
financial NLP applications the natural unit of analysis is longer: an
earnings call transcript (several thousand words), a 10-K filing
(tens of thousands of words), a news headline (ten words), or a
research note (hundreds of words).  To use embedding-based features
at the document level we need an aggregation scheme that maps a
sequence of token embeddings to a single fixed-size vector.

% ----------------------------------------------------------------
\subsection{Averaging Word Vectors}
\label{sec:doc-repr-avg}

The simplest aggregation is the unweighted arithmetic mean.

\begin{definition}[Mean Embedding]
\label{def:mean-embedding}
Let $d = (w_1, w_2, \dots, w_n)$ be a document of $n$ tokens and let
$\varphi : \mathcal{V} \to \mathbb{R}^d$ be a word embedding.
The \emph{mean embedding} of $d$ is
\begin{equation}
  \bar{\boldsymbol{\varphi}}(d)
  \;=\;
  \frac{1}{n} \sum_{i=1}^{n} \varphi(w_i)
  \;\in\; \mathbb{R}^d.
  \label{eq:mean-embedding}
\end{equation}
\end{definition}

\paragraph{Connection to Bag-of-Words.}
Mean embedding is algebraically equivalent to multiplying the
Bag-of-Words term-frequency vector $\mathbf{tf}(d) \in \mathbb{R}^V$
by the (column-normalised) embedding matrix.  Specifically, let
$\mathbf{tf}(d)_j = \text{count}(w_j, d)/n$ be the relative frequency
of word $w_j$ in document $d$.  Then
\begin{equation}
  \bar{\boldsymbol{\varphi}}(d)
  \;=\;
  \mathbf{W}^{\top} \mathbf{tf}(d)
  \;\in\; \mathbb{R}^d,
  \label{eq:mean-embedding-bow}
\end{equation}
where $\mathbf{W} \in \mathbb{R}^{V \times d}$ is the embedding matrix
from Definition~\ref{def:word-embedding}.  This shows that mean
embedding is a linear operation on the standard BoW representation,
projecting from $V$-dimensional one-hot/count space into the dense
$d$-dimensional embedding space via the weight matrix $\mathbf{W}$.

\begin{proposition}
\label{thm:mean-embed-is-bow}
Mean embedding is a Bag-of-Words model.  Specifically, if
$\hat{\mathbf{W}} = [\varphi(w_1) \;\cdots\; \varphi(w_V)]^{\top}
\in \mathbb{R}^{V \times d}$ is the full embedding matrix, then
\[
  \bar{\boldsymbol{\varphi}}(d)
  \;=\;
  \hat{\mathbf{W}}^{\top} \mathbf{tf}(d).
\]
In particular, the mean embedding is invariant to permutations of the
token sequence; word order is discarded.
\end{proposition}
\begin{proof}
  Expand \eqref{eq:mean-embedding} using the row-selection
  interpretation $\varphi(w_i) = \hat{\mathbf{W}}^{\top}\mathbf{e}_{w_i}$,
  then use linearity:
  \[
    \frac{1}{n}\sum_{i=1}^n \hat{\mathbf{W}}^{\top}\mathbf{e}_{w_i}
    \;=\;
    \hat{\mathbf{W}}^{\top}
    \!\underbrace{\left(\frac{1}{n}\sum_{i=1}^n \mathbf{e}_{w_i}\right)}_{\mathbf{tf}(d)}
    \;=\;
    \hat{\mathbf{W}}^{\top}\mathbf{tf}(d). \qedhere
  \]
\end{proof}

\paragraph{Limitations of mean embedding.}
Despite its simplicity and computational efficiency, mean embedding
has several well-documented weaknesses:
\begin{enumerate}
  \item \textbf{Equal weighting.}  Common words such as \emph{the},
        \emph{is}, and \emph{of} contribute as much to the mean as
        domain-specific terms such as \emph{default} or
        \emph{liquidity}.  Stop-word removal partially addresses this,
        but introduces a pre-processing dependency.
  \item \textbf{No word order.}  The mean embedding of
        ``\emph{the company did not meet earnings expectations}'' is
        identical to that of ``\emph{earnings expectations met the
        company not did}'' — a garbled sentence that contains the same
        words.  Syntactic structure, negation, and conditionals are
        invisible.
  \item \textbf{Poor performance on long documents.}  As $n$ grows,
        the mean converges (by the law of large numbers) toward the
        corpus-wide average embedding, losing document-specific
        signal.  Empirically, mean embedding performs well for short
        texts (headlines, short paragraphs) but degrades for filings
        or call transcripts \cite{reimers2019}. % [CITE: Reimers & Gurevych (2019). Sentence-BERT: Sentence Embeddings using Siamese BERT-Networks. EMNLP]
  \item \textbf{Out-of-vocabulary words.}  Words absent from
        $\mathcal{V}$ (e.g.\ rare tickers, neologisms) must be mapped
        to a special \texttt{<UNK>} vector or dropped, both of which
        introduce noise.
\end{enumerate}

% ----------------------------------------------------------------
\subsection{Weighted Aggregation and Sentence Transformers}
\label{sec:doc-repr-sbert}

\subsubsection{TF-IDF Weighted Averaging}

A natural refinement is to weight each word's embedding by its
TF-IDF score, up-weighting terms that are rare in the corpus but
frequent in the document and down-weighting common uninformative words.

\begin{definition}[TF-IDF Weighted Embedding]
\label{def:tfidf-embedding}
Let $\text{tfidf}(w_i, d)$ denote the TF-IDF weight of token $w_i$
in document $d$ (see Section~\ref{sec:classical-tfidf}).  The
\emph{TF-IDF weighted embedding} is
\begin{equation}
  \boldsymbol{\varphi}_{\mathrm{tfidf}}(d)
  \;=\;
  \frac{\displaystyle\sum_{i=1}^{n} \text{tfidf}(w_i, d)\; \varphi(w_i)}
       {\displaystyle\sum_{i=1}^{n} \text{tfidf}(w_i, d)}.
  \label{eq:tfidf-embedding}
\end{equation}
\end{definition}

This is a weighted arithmetic mean with weights $\alpha_i =
\text{tfidf}(w_i, d) / \sum_j \text{tfidf}(w_j, d) > 0$.  By
design, $\sum_i \alpha_i = 1$, so $\boldsymbol{\varphi}_{\mathrm{tfidf}}(d)$
lies in the convex hull of the token embeddings.  TF-IDF weighting
consistently outperforms uniform averaging on document classification
and similarity tasks when the embedding quality is held fixed
\cite{arora2017}. % [CITE: Arora et al. (2017). A Simple but Tough-to-Beat Baseline for Sentence Embeddings. ICLR]

\begin{remark}
\label{rem:sif}
A closely related approach, \emph{Smooth Inverse Frequency} (SIF)
\cite{arora2017}, % [CITE: same as above]
replaces TF-IDF weights with
$\alpha_i = a / (a + p(w_i))$,
where $p(w_i)$ is the unigram probability of word $w_i$ in the
training corpus and $a \approx 10^{-3}$ is a smoothing constant.
After computing the weighted sum, SIF removes the first principal
component, which typically captures uninformative ``common discourse''
variance.  SIF embeddings match or exceed supervised methods on
several semantic textual similarity benchmarks at negligible
computational cost.
\end{remark}

\subsubsection{The Limitation of Static Embeddings: Polysemy}

Both Word2Vec and GloVe assign a \emph{single, fixed} vector to each
vocabulary item.  This is adequate for words with a single dominant
sense, but fails for polysemous words — words with multiple distinct
meanings depending on context.  In finance, this problem is pervasive:
\begin{itemize}
  \item \emph{Bank} can mean a financial institution or the riverbank.
  \item \emph{Yield} can mean bond yield, crop yield, or to give way.
  \item \emph{Position} can mean a financial position (long/short) or
        a physical location.
  \item \emph{Default} can mean failure to repay a debt or the
        default setting of a software parameter.
\end{itemize}
Under a static embedding, the vector for \emph{bank} is a single
point in $\mathbb{R}^d$ that is a weighted average of both senses,
equidistant from financial and geographic contexts.  This conflation
can harm downstream tasks: a sentiment model that associates
\emph{bank} with positive sentiment in financial reporting will
incorrectly transfer that sentiment to sentences about river flooding.

\subsubsection{Contextual Embeddings and Sentence-BERT}

\paragraph{From static to contextual embeddings.}
The solution is to produce embeddings that depend on the full
surrounding context.  ELMo \cite{peters2018} % [CITE: Peters et al. (2018). Deep contextualized word representations. NAACL]
introduced deep contextual embeddings from a bidirectional language
model; BERT \cite{devlin2019} % [CITE: Devlin et al. (2019). BERT: Pre-training of Deep Bidirectional Transformers for Language Understanding. NAACL]
later generalised this with a transformer architecture.  In a BERT
encoder, the representation of the token \emph{bank} at layer $\ell$
is a function of the entire input sentence, so it will differ in
``\emph{The Federal Reserve Bank raised rates}'' versus
``\emph{The river bank flooded after the storm}.''

\paragraph{Sentence-BERT (SBERT).}
Vanilla BERT was designed for token-level tasks (named entity
recognition, question answering) rather than for producing
sentence- or document-level representations.  Naively averaging the
BERT token representations (mean pooling) or using the \texttt{[CLS]}
token produces suboptimal sentence embeddings for semantic similarity
tasks \cite{reimers2019}. % [CITE: Reimers & Gurevych (2019). Sentence-BERT. EMNLP]

Sentence-BERT (SBERT) \cite{reimers2019} % [CITE: same]
addresses this by fine-tuning a BERT encoder on a Natural Language
Inference (NLI) objective using a \emph{Siamese} architecture:
two weight-sharing BERT encoders process a sentence pair
$(s_a, s_b)$, producing pooled representations
$\mathbf{u} = \text{POOL}(\text{BERT}(s_a))$ and
$\mathbf{v} = \text{POOL}(\text{BERT}(s_b))$.
The model is trained to predict whether $s_a$ and $s_b$ are
semantically equivalent (entailment), contradictory, or unrelated,
using a classification head over the concatenation
$[\mathbf{u};\, \mathbf{v};\, |\mathbf{u} - \mathbf{v}|]$.

\begin{definition}[Mean Pooling in SBERT]
\label{def:sbert-pool}
Let $\mathbf{H} = [\mathbf{h}_1, \mathbf{h}_2, \dots, \mathbf{h}_n]
\in \mathbb{R}^{n \times d_\text{model}}$
be the matrix of last-layer hidden states from the BERT encoder for an
$n$-token input sequence (including special tokens \texttt{[CLS]} and
\texttt{[SEP]}).  Under \emph{mean pooling}, the sentence embedding is
\begin{equation}
  \mathbf{s}
  \;=\;
  \frac{1}{n}\sum_{i=1}^{n} \mathbf{h}_i \cdot m_i,
  \label{eq:sbert-mean-pool}
\end{equation}
where $m_i \in \{0, 1\}$ is an attention mask excluding padding tokens.
\end{definition}

\paragraph{Mean pooling vs.\ \texttt{[CLS]} token.}
BERT's \texttt{[CLS]} token is prepended to every input sequence and
its final hidden state $\mathbf{h}_{\texttt{[CLS]}}$ is used in
classification tasks; one might expect it to encode sentence-level
meaning.  In practice, without explicit fine-tuning on semantic
similarity tasks, the \texttt{[CLS]} representation performs
surprisingly poorly as a general-purpose sentence embedding
\cite{reimers2019}. % [CITE: same]
Mean pooling over all token representations consistently outperforms
\texttt{[CLS]} pooling on standard semantic textual similarity (STS)
benchmarks.  After SBERT fine-tuning, both pooling strategies converge
to high-quality representations, with mean pooling retaining a slight
advantage.

\paragraph{Computational efficiency.}
A key practical advantage of SBERT is that sentence embeddings can be
\emph{pre-computed and indexed}: once the corpus of $N$ documents has
been encoded, semantic search reduces to a nearest-neighbour query in
$\mathbb{R}^{d_{\text{model}}}$, which can be accelerated with
approximate nearest-neighbour (ANN) data structures (e.g.\ FAISS
\cite{johnson2019}, % [CITE: Johnson et al. (2019). Billion-scale similarity search with GPUs. IEEE TPAMI]
HNSW).  Cross-encoding a query against all $N$ documents with vanilla
BERT would require $N$ separate forward passes; SBERT bi-encoding
requires only one forward pass per document, reducing inference time
from $O(N)$ to $O(1)$ at query time.

\paragraph{Financial applications.}

\begin{example}
\label{ex:doc-similarity-finance}
The following applications of contextual document embeddings are
documented in the finance literature:
\begin{enumerate}
  \item \textbf{Earnings call similarity.}
        Encode all earnings call transcripts in a corpus using SBERT
        (or a finance-tuned variant such as FinBERT-SBERT).  At each
        reporting period, retrieve the $k$ most similar past calls
        using cosine similarity.  The similarity score can serve as a
        feature for predicting short-window abnormal returns
        \cite{ko2020}. % [CITE: Ko et al. (2020). or relevant FinNLP earnings similarity paper]
  \item \textbf{Merger target identification.}
        Encode 10-K business description sections for all public firms.
        High cosine similarity between two firms' descriptions predicts
        product-market proximity, which is a strong predictor of
        acquisition likelihood \cite{frattaroli2019}. % [CITE: relevant M&A text paper]
  \item \textbf{Regulatory document search.}
        A compliance team encodes a large corpus of regulatory guidance
        documents (Basel III, IFRS 9, MiFID II).  When a new
        regulation is published, SBERT-based retrieval identifies
        the most relevant precedents, replacing keyword search with
        semantic search.
  \item \textbf{Factor model narrative alignment.}
        Embed analyst research notes and quarterly earnings guidance.
        Measure the cosine similarity between the guidance embedding
        and subsequent actual results narratives as a proxy for
        management credibility or guidance quality.
\end{enumerate}
\end{example}

\begin{remark}
\label{rem:context-window-limit}
BERT-based models impose a maximum context length (typically 512
tokens for standard BERT, 4096 or more for long-context variants such
as Longformer \cite{beltagy2020}). % [CITE: Beltagy et al. (2020). Longformer: The Long-Document Transformer. arXiv:2004.05150]
Long financial documents (full 10-K filings, lengthy earnings call
transcripts) must be handled by one of three strategies:
(i) \emph{truncation} — only the first 512 tokens are encoded, which
may miss material information in later sections;
(ii) \emph{chunking} — the document is split into overlapping windows,
each encoded independently, and the chunk embeddings are aggregated
(e.g.\ by mean pooling or weighted by position); or
(iii) \emph{hierarchical encoding} — sentences are encoded individually
and a second-level encoder (another transformer or an LSTM) aggregates
sentence embeddings into a document embedding.  The choice among these
strategies introduces a design trade-off between completeness and
computational cost that has no universal solution; practitioners should
benchmark all three on their specific downstream task.
\end{remark}

% ---------------------------------------------------------------
\section{Sequential Models: From RNNs to Attention}
\label{sec:sequential}

\subsection{Recurrent Neural Networks and the Vanishing Gradient Problem}
\label{sec:sequential-rnn}

[Placeholder]

\subsection{LSTMs and GRUs}
\label{sec:sequential-lstm}

[Placeholder]

\subsection{The Attention Mechanism}
\label{sec:sequential-attention}

[Placeholder]

% ---------------------------------------------------------------
\section{The Transformer Architecture}
\label{sec:transformer}

\subsection{Self-Attention (Mathematics)}
\label{sec:transformer-selfattn}

[Placeholder]

\subsection{Multi-Head Attention}
\label{sec:transformer-mha}

[Placeholder]

\subsection{Positional Encoding}
\label{sec:transformer-pos}

[Placeholder]

\subsection{Encoder--Decoder Structure}
\label{sec:transformer-encdec}

[Placeholder]

\subsection{Pre-training, Fine-tuning, and BERT vs.\ GPT}
\label{sec:transformer-pretrain}

[Placeholder]

% ---------------------------------------------------------------
\section{The Modern LLM Landscape}
\label{sec:landscape}

\subsection{Generative Models vs.\ Reasoning Models}
\label{sec:landscape-genreason}

[Placeholder]

\subsection{Major Model Families}
\label{sec:landscape-families}

[Placeholder]

\subsection{Financial NLP Benchmarks and Datasets}
\label{sec:landscape-benchmarks}

[Placeholder]

% ---------------------------------------------------------------
\section{Working with LLMs via API}
\label{sec:api}

\subsection{Core Concepts: Tokens, Context Windows, and Sampling}
\label{sec:api-concepts}

[Placeholder]

\subsection{Python Walkthrough: OpenAI, Anthropic, and HuggingFace}
\label{sec:api-python}

[Placeholder]

\subsection{Cost, Latency, and Practical Considerations}
\label{sec:api-practical}

[Placeholder]

% ---------------------------------------------------------------
\section{Limitations and Responsible Use}
\label{sec:limitations}

\subsection{Hallucinations in Financial Contexts}
\label{sec:limitations-hallucinations}

[Placeholder]

\subsection{Regulatory and Ethical Considerations}
\label{sec:limitations-ethics}

[Placeholder]

% ---------------------------------------------------------------
\section{Roadmap of This Book}
\label{sec:roadmap}

[Placeholder]
